% =====================================================================
\section{Study design}
\label{sec:design}
% Source: docs/paper/PREREGISTRATION.md (sections 1-9; amendments 0.1 A1-A9).
% =====================================================================

The study is a pre-registered fault-injection experiment. We take clean stage
outputs from real pipeline runs, corrupt each with one of six operators, and
present the corrupted output to the production agent for that stage under
controlled variations of what the supervisor says and where the defective
output appears. The pre-registration fixes the estimand, the arms and their
wording, the outcome coding, the exclusions, the single confirmatory test, the
power target, and the reading of every exploratory cell. It was frozen as the
annotated git tag \code{prereg-freeze} on 14~September~2026. The first trial
ran on 5~October~2026. Every later change is a dated amendment, and
Section~\ref{sec:design-deviations} lists them all.

% ---------------------------------------------------------------------
\subsection{Estimand and research questions}
\label{sec:design-rq}

The pre-registered estimand is a single question. \emph{For a defect injected
into a clean stage output and presented to the production agent as an artifact
under review, does the content of the supervisor's feedback change whether the
regenerated output is correct?} Three things are excluded from it by design.
The unit is the defect, not the run, so ``repaired'' does not mean that the run
became good. The retried stage regenerates its whole output rather than editing
the faulty item, so the construct is \emph{repair by regeneration}. And the
injected defects are not a sample of production defects, so no repair rate in
this paper is a production repair rate. The first exclusion is revisited by
RQ4 and the third by RQ3; both are exploratory.

\begin{description}[leftmargin=1.2em,style=nextline]
\item[RQ1 (confirmatory).] Does the content of the supervisor's feedback (a
  measured fact versus a contentless rejection) change whether a retried stage
  repairs an injected defect?
\item[RQ2.] Why does repair succeed or fail? We vary the placement of the
  defective artifact, the sampling temperature, the extent of the defect, and
  whether the stage's obligations are restated at retry time.
\item[RQ3.] What does the production retry do on real, production-caught
  defects, and how does that compare with injected ones?
\item[RQ4.] Does stage-level repair become pipeline-level improvement?
\item[RQ5.] What does deterministic recompute-and-overwrite buy over a retry?
\item[RQ6.] Do the answers hold for a second model from a different family?
  \CSEVEN{RQ6 is answered by the second-model cell (plan D5b; amendment A10 to
  be filed before it runs).}
\end{description}

Only RQ1 carries a hypothesis test. Everything else is an estimate with an
interval, labelled exploratory or descriptive where it appears.

% ---------------------------------------------------------------------
\subsection{Workloads and injection bases}
\label{sec:design-workloads}

The workloads are restaurants from the Yelp Open Dataset~\citep{yelpdataset}.
The sampling pool is the 52{,}268 businesses in the \emph{Restaurants} category
with at least five reviews. Review count per business is heavily skewed
(p25\,=\,14, p50\,=\,34, p75\,=\,92, p90\,=\,214), so the sample is stratified
at those four percentiles rather than drawn uniformly. Under a fixed seed, the
sampler takes the seven businesses whose review counts lie nearest each target,
giving 28 workloads. The first twenty, five per stratum, are the workloads of
the conference study~\citep{hoang2026eas}. The other eight were drawn by the
same sampler before the freeze, as a strict extension of the first twenty, to
reach the power target of Section~\ref{sec:design-analysis}. No business was
added or removed by hand. Each run loads a business's complete review set.

Each workload was run once through the pipeline in \emph{observe mode}, in
which the supervisor measures every stage and records its verdict but always
proceeds. The model was \model{} at temperature~0. Every stage's input and
output was saved. These run records are the injection bases. A stage output is
eligible as a base only if the supervisor's own measurement of it was clean,
so that the injected defect is the only defect at that stage.

% ---------------------------------------------------------------------
\subsection{Fault operators}
\label{sec:design-operators}

Table~\ref{tab:operators} lists the operators. Each corrupts one clean stage
output so as to trip exactly one class of supervisor check. Every random
choice is drawn from a stream seeded by business, operator and seed, so each
injected output is reproducible byte for byte. The operator and the supervisor
are independent programs. The operator records what it did and which check it
expects to fire; the supervisor's \code{measure()} reports what actually fired.

\begin{table}[tbp]
\centering
\caption{Fault operators. ``Core'' is the number of defects in the
confirmatory core after both validity gates. \op{O4} and \op{O8} are not repair
operators. \op{O4} produces a warning rather than a retry, because production
overwrites the frequency (Section~\ref{sec:system-overwrite}); it runs only in
the overwrite counterfactual. \op{O8} is a positive control whose repair is
guaranteed, because the report prompt carries the correct values; it was not
run.}
\label{tab:operators}
\footnotesize
\begin{tabularx}{\linewidth}{@{}l l L{2.6cm} Y r@{}}
\toprule
Op. & Stage & Check targeted & What it does & Core \\
\midrule
\op{O1} & reasoning & unresolvable evidence id & deletes one character from one
  real evidence id of one pattern & 24 \\
\op{O2} & reasoning & unresolvable evidence id & appends a well-formed id that
  no analysed review carries & 24 \\
\op{O3} & reasoning & misaligned aspect & relabels a pattern with an aspect
  that none of its evidence reviews mentions & 22 \\
\op{O4} & reasoning & frequency divergence (warning) & moves a pattern's
  frequency 0.15 away from the recomputed value & (23) \\
\op{O5} & strategy & untraceable recommendation & replaces one
  recommendation's issue with an issue from a fixed corpus that traces to
  nothing upstream & 20 \\
\op{O6} & report & report containment & inserts a root cause produced by no
  upstream stage & 26 \\
\op{O7} & report & report containment & inserts a recommendation produced by
  no upstream stage & 26 \\
\op{O8} & report & trusted metadata & perturbs the business name or the
  sample size & -- \\
\midrule
\multicolumn{4}{@{}l}{Repair operators (\op{O1}--\op{O3}, \op{O5}--\op{O7})} & 142 \\
\bottomrule
\end{tabularx}
\end{table}

The six repair operators over 28 bases give 168 candidate defects. For 26 of
them the operator had nothing to corrupt or the base was not clean, which
leaves 142. Two validity gates then run on each injected output, before any
model call. The first, \emph{injection verified}, requires that the targeted
check fires and that no other retry-level check fires that did not fire on the
base. The second, \emph{triggered}, requires that the supervisor's verdict on
the injected output is \verdict{retry}. A defect that is not triggered receives
no feedback of any kind, so the arms would be identical for it; such defects
are removed and counted rather than scored as unrepaired, following
\citet{tan2026agentchaos}. Both gates passed 142 of 142 defects, and the result
was stable across four seeds. This should not be read as discrimination: the
applicability rules already encode most of what the gates check. The gates are
a standing guarantee, not a filter on this bank.

\paragraph{How the injected defects compare with production defects}
Before the study, two raters independently judged a blinded sheet of 31
defects: 11 that the supervisor had caught in production and 20 injected
ones. The sheet hid each defect's origin. Agreement was 29 of 31, with Cohen's
$\kappa = 0.870$ (95\% CI \ci{0.674}{1.000}). The raters called none of the 11
production-caught defects an error; they judged all of them technicalities
or checker false positives. They called 16 of the 20 injected defects errors.
The split followed mechanism, not origin. \op{O1}, which corrupts a
\emph{real} identifier, was judged a technicality in all four of its items,
like the production defects. Fabrication (\op{O2}), mislabelling (\op{O3}) and
foreign content (\op{O5}--\op{O7}) were judged errors in all 16 of their items.
This audit is why the paper never pools absolute repair rates across
operators. \op{O1} is reported as the production-like cell and
\op{O2}--\op{O7} as the severe cell. The contrast between arms is still pooled,
for the reason given in Section~\ref{sec:design-analysis}.

% ---------------------------------------------------------------------
\subsection{Arms and placements}
\label{sec:design-arms}

A trial re-runs the defective stage once through the production agent. Two
factors vary: what the model is told (the \emph{arm}) and where the defective
output appears in its context (the \emph{placement}). Both are built on top of
the production prompt for the stage, so arms differ only in the turns they add.

\begin{description}[leftmargin=1.2em,style=nextline]
\item[\arm{A1}, measured-fact feedback.] The production feedback template
  (Section~\ref{sec:system-feedback}), filled with the measured fact that
  triggered the verdict. This is what the production supervisor sends.
\item[\arm{A2}, contentless rejection.] The same template with the measured
  facts removed: an authoritative rejection that says nothing about what is
  wrong.
\item[\arm{A0}, capability ceiling.] The defective output is presented as a
  draft to review, with no retry framing. It bounds what any feedback can
  achieve.
\item[\arm{A3}, regeneration baseline.] The original stage prompt and nothing
  else.
\end{description}

\noindent The placement fixes how, and whether, the defective output enters
the context:

\begin{description}[leftmargin=1.2em,style=nextline]
\item[\arm{P0}, artifact absent.] The model does not see the defective output.
  This is the production path, because the production retry appends only the
  feedback to the stage prompt.
\item[\arm{P1}, artifact under review.] The defective output appears in a user
  turn as the previous answer to the task.
\item[\arm{P2}, artifact as the model's own turn.] The defective output appears
  as a prior assistant turn.
\item[\arm{P3}, own turn with a marker.] As \arm{P2}, followed by a user turn
  stating that the answer above contains an error. This level is exploratory.
\end{description}

\noindent The design is not fully crossed, and two cells are undefined by
construction. \arm{A0} needs an artifact to correct, so it does not exist at
\arm{P0}. \arm{A3} stops being a bare regeneration once it is shown an
artifact, so it exists only at \arm{P0}. \ref{app:wording} gives the exact
wording of every arm and placement; it was frozen in code on
25~August~2026, before the pre-registration and before any trial.

\paragraph{Why the core runs at \arm{P1}}
At \arm{P0} an injected defect never enters the model's context. A retry there
would be dominated by the model reproducing its own clean output, and the
difference between arms would measure only how feedback text perturbs
generation. The confirmatory core therefore runs \arm{A1} and \arm{A2} at
\arm{P1}, and \arm{P0} becomes a diagnostic level of the placement cell. The
replay of real defects runs at \arm{P0}, because there the defect genuinely was
the model's own and the production retry genuinely does not show it. One
consequence follows for \arm{A3}. It exists only at \arm{P0}, so on injected
defects it regenerates from a clean context and is \emph{not} a floor. It is a
regeneration baseline: what a fresh draw that never saw the defect scores.
This was established before the freeze, and A3 was re-roled on that basis
(Section~\ref{sec:design-deviations}).

% ---------------------------------------------------------------------
\subsection{Outcomes and scoring}
\label{sec:design-outcomes}

Each trial's regenerated output is scored into one of four mutually exclusive
classes, decided in a fixed order:
\begin{enumerate}[label=(\roman*)]
\item the injected check still fires: \cls{unrepaired};
\item the targeted item is gone, for an operator whose clean output had it:
  \cls{deleted};
\item the item is back to its pre-injection value: \cls{repaired};
\item otherwise: \cls{replaced}.
\end{enumerate}
The confirmatory analysis uses the strict dichotomy of \cls{repaired} versus
everything else. Two scoring decisions matter for reading the results. First,
success is the \emph{replacement value}, not the disappearance of the
complaint. A fabricated identifier swapped for a different identifier that
happens to resolve passes the check, but it scores \cls{replaced}, because the
output is still wrong. Second, \cls{deleted} counts against repair. Dropping
the item the supervisor complained about is evasion, and scoring it as success
would let ``remove the pattern you were asked about'' register as a supervisor
win. For the insertion operators \op{O6} and \op{O7}, removing the invented
item \emph{is} the repair, so those trials score \cls{repaired} and
\cls{deleted} cannot occur. Items are matched by identity, not position,
because a regenerated output reorders its items. Each substitution operator
anchors on a field it does not touch.

\paragraph{Scoring \op{O5} on content words}
The production traceability check compares a recommendation with upstream
items by a character-level similarity ratio, with a threshold of 0.35. On
pre-study observe-mode data, that check flagged 8 of 49 correct strategy
outputs. Scoring \op{O5} with it would therefore have scored correct
regenerations as \cls{unrepaired} at a similar rate in every arm. The scorer
decides whether the \op{O5} check still fires by content-word Jaccard
similarity of at least 0.10 against the same upstream candidates. The identity
anchor (the recommendation's \code{action}) is matched at a Jaccard threshold
of 0.20. Both thresholds were fixed on calibration data before any trial
outcome existed. The production supervisor itself is unchanged: the \op{O5}
bases, the injected issue, the trigger and the \arm{A1} feedback are identical
under both matchers. The production verdict is logged beside the scorer's and
reported as a sensitivity analysis, together with containment\,$>0$ and
Jaccard\,$\ge 0.05$.

\paragraph{Collateral damage}
Collateral damage is a second dimension, never a class. It is the loss of
\emph{non-injected} items relative to the pre-injection output, counted above
the churn that an unperturbed temperature-0 pipeline already produces. That
churn was measured in a repetition study before this one and differs by item
class (0.07 for pattern aspects, 0.33 for pattern descriptions, 0.29 for root
causes, 0.23 for recommendation issues). The injected item is excluded, so the
outcome class and the collateral flag are never raised by the same event.
Report-stage items have no calibrated floor, so those trials are excluded from
the collateral denominator rather than counted as undamaged.

% ---------------------------------------------------------------------
\subsection{Analysis}
\label{sec:design-analysis}

\paragraph{Unit and estimate}
The unit of analysis is the defect. Within each (defect, arm) cell the $r=3$
repetitions are averaged into one rate. The primary estimate is the mean over
defects of the per-defect difference \arm{A1}\,$-$\,\arm{A2}. Its 95\%
interval comes from a bootstrap that resamples whole defects (10{,}000
resamples, fixed seed), which keeps repetitions of one defect from counting as
independent evidence. A defect enters the contrast only if it has at least one
usable repetition in both arms. The estimate, not the test, is the primary
reported quantity, so that a null result bounds the effect rather than merely
failing to reject.

\paragraph{The confirmatory test}
There is exactly one: an exact McNemar test of \arm{A1} against \arm{A2} at
$\alpha=0.05$, on a per-defect dichotomy. A defect counts as repaired in an arm
if strictly more than half of its completed repetitions were repaired. Two
alternatives were rejected before the freeze. \emph{Any-of-three} measures
best-of-$k$ and rewards variance. \emph{Unanimity} penalises variance and would
make the test weaker than the estimate it accompanies. A one-to-one split
after a lost repetition resolves to not repaired, which cannot bias a matched
contrast.

\paragraph{Power and the sample size}
Power was computed by simulation, because repetitions nest within defects and
all arms of one defect share a serving session. At the delivered $N=142$ and
$r=3$, power was 0.875 for a 12\pp{} effect, and the minimum detectable effect
was 10.9\pp{}, with the shared-session (``occasion'') effect included. The
sample was grown to reach this. At the 102 defects the first twenty workloads
yield, the detectable effect was 12.6\pp{}, above the 12\pp{} threshold the
study had set for itself. Rather than lower the threshold, eight more workloads
were collected.

\paragraph{Pooling the contrast, not the rates}
The primary contrast is estimated over all six repair operators, although the
audit separates \op{O1} from the rest. The audit bounds the external validity
of an absolute repair rate. Whether feedback content changes the outcome is a
different claim, and neither literature it adjudicates frames that claim as
holding only within a severity band. To expose the pooling assumption, the
contrast is also reported within \op{O1} ($N=24$) and within the severe cell,
unconditionally and in the same table.

\paragraph{Pre-specified secondary analyses}
These analyses are named in advance, with their interpretation fixed, but they
carry no correction for multiplicity and support no headline claim:
\begin{itemize}
\item the \op{O1}\,$\times$\,arm interaction;
\item the \arm{A1}/\arm{A0} ratio;
\item \arm{A3} by operator;
\item per-class repair rates against a recorded directional prediction;
\item stratification by defect origin (real versus injected) and by stage;
\item the chain ``did the output change'' $\rightarrow$ ``did the defect
  resolve'';
\item the terminal supervisor verdict per trial;
\item \op{O5} sensitivity, including the primary contrast without \op{O5};
\item \arm{A2} against \arm{A3}, with Cochran's $Q$ across the three arms they
  share, read as re-presenting the defect with a rejection versus regenerating
  without it.
\end{itemize}

\paragraph{Exclusions}
Every exclusion is counted. A defect that fails a validity gate is removed
before any model call, so it leaves all arms at once and cannot break the
matching. A trial whose call ends in a provider error (\code{status\_error}) has
no interpretation under any arm and is excluded. Such a trial is never re-run
until it succeeds, because that would resample on the outcome. A trial served
by the fallback model would also be excluded. Fallback was disabled, so a
provider failure surfaces as an error rather than as a silent change of model.
A trial whose first call fails schema validation is re-issued with the
model's own invalid output in an assistant turn, which embeds a
\arm{P2}-like turn in a \arm{P0} or \arm{P1} trial. Such trials are logged and
analysed separately, not excluded, and the primary estimate is reported both
with and without them.

\paragraph{The publication rule}
One further rule was fixed before the sample was sized: the repair study would
become a paper of its own if and only if the \arm{A1}\,$-$\,\arm{A2} estimate
were at least 12\pp{} and its interval excluded zero. Neither condition was
met (Section~\ref{sec:results-rq1}), so the study is reported here as an
extension of the system paper.

% ---------------------------------------------------------------------
\subsection{Exploratory cells}
\label{sec:design-cells}

\begin{table}[tbp]
\centering
\caption{Every cell that ran, with its role. The confirmatory core is the only
cell frozen with the pre-registration. Each other cell was specified in a
dated amendment, committed before it ran
(Section~\ref{sec:design-deviations}). ``Usable'' excludes provider errors.
Cells marked $\dagger$ re-run their own in-session comparators instead of
borrowing the core's (Section~\ref{sec:design-execution}).}
\label{tab:cells}
\footnotesize
\begin{tabularx}{\linewidth}{@{}L{2.8cm} l Y r r r@{}}
\toprule
Cell & Amd. & Design & Defects & Trials & Usable \\
\midrule
Core (confirmatory) & -- & \arm{A1}, \arm{A2} at \arm{P1} (all defects);
  \arm{A0} at \arm{P1} (half); \arm{A3} at \arm{P0} (a quarter); $r=3$ & 142 &
  1{,}167 & 1{,}149 \\
Overwrite counterfactual & A2 & \op{O4}; \arm{A1}@\arm{P1}, \arm{A3}@\arm{P0}
  & 23 & 138 & 138 \\
Constraint disclosure$^\dagger$ & A1, A3 & \{C0, C1\} $\times$ \{\arm{A1},
  \arm{A2}\} at \arm{P1} & 37 & 444 & 444 \\
Placement$^\dagger$ & A5 & \arm{A1} $\times$ \{\arm{P0}--\arm{P3}\} & 45 & 540
  & 539 \\
Severity$^\dagger$ & A6 & \op{O1} levels 1, 2; \arm{A1}@\arm{P1} & 24 & 144 &
  144 \\
Warm resampling$^\dagger$ & A7 & \arm{A1}, \arm{A2}@\arm{P1} at $T=0$ and $T=1$;
  \arm{A3}@\arm{P0} at $T=1$ & 45 & 675 & 675 \\
Real-defect replay & A8 & 13 production-caught defects; \arm{A1}--\arm{A3}
  at \arm{P0}; descriptive & 13 & 117 & 117 \\
Downstream propagation & A9 & later stages re-run from clean, unrepaired,
  \arm{A1}, \arm{A2} outputs & 90 & 447$^\ast$ & 447 \\
\CSEVEN{Second model} & \CSEVEN{A10} & \CSEVEN{core contrast on a second model}
  & \CSEVENNUM & \CSEVENNUM & \CSEVENNUM \\
\bottomrule
\multicolumn{6}{@{}l}{$^\ast$ Conditions (defect $\times$ starting point), 796
calls in total.}
\end{tabularx}
\end{table}

Table~\ref{tab:cells} lists the cells. Each exploratory cell answers one part
of RQ2--RQ5.

\begin{description}[leftmargin=1.2em,style=nextline]
\item[Constraint disclosure (C1).] Before the freeze, an audit of the prompts
  found that the \arm{A1} feedback asserts five obligations that no stage prompt
  states. This is the prompt-completeness confound that \citet{huang2024cannot}
  and \citet{kamoi2024} warn of: A1 might win by teaching the task rather than
  by naming the defect. Condition C1 splices the stage's own obligations into
  the retry prompt; C0 is the production prompt. Both conditions inject into the
  same base, so the cell measures \emph{stating the constraints at retry time},
  not a constraint-complete pipeline. Detecting the interaction had power of
  0.11--0.30 at sample sizes out of reach, so the cell was pre-specified as two
  gaps with intervals, never a test.
\item[Placement.] \arm{A1} at all four placements, on a third of the defects
  selected by a stratified per-operator stride. The pre-recorded prediction was
  \arm{P1}\,$\gg$\,\arm{P3}\,$\approx$\,\arm{P2}, with \arm{P0} reported apart.
\item[Severity.] \op{O1} at two levels. Level~1 is the frozen operator. Level~2
  corrupts \emph{every} resolving evidence identifier of the targeted pattern
  by one character, so per-identifier damage is unchanged and only its extent
  varies. This definition was chosen before either level ran, from two
  candidates defined together.
\item[Warm resampling.] Temperature~1.0, to test whether the ceiling is an
  artefact of greedy decoding and to give the equal-budget comparison with
  best-of-$k$ resampling that \citet{olausson2024selfrepair} require. This cell
  was designed \emph{after} the core's ceiling and its byte-identity pattern
  were known. Both possible readings were fixed before it ran.
\item[Real-defect replay.] The 13 usable defects that the production
  supervisor caught in the 28 observe-mode runs, in 7 business\,$\times$\,stage
  cells. Eight are untraceable recommendations at the strategy stage. At the
  reasoning stage, three are unresolvable evidence identifiers and two are
  misaligned aspects. A defect runs only if replaying its stored run reproduces
  the recorded retry exactly (verdict, fired checks and feedback string); all
  13 do. There is no clean reference for a real defect, so the outcome is
  whether the supervisor \emph{accepts} the retried output. Intervals resample
  cells, not defects, and no test is run. At this scale a cluster-bootstrap
  test over-rejects: by simulation, its type~I error with five clusters was
  0.086 at a nominal 0.05.
\item[Downstream propagation.] For the 90 core defects with a later stage
  (\op{O1}--\op{O3}, \op{O5}), the later stages are re-run once each through
  the production nodes, in observe mode, from five starting points: the clean
  output twice (the second run is the in-session noise floor), the unrepaired
  injected output, and the core's \arm{A1} and \arm{A2} retried outputs. Three
  endpoints were fixed before the pass. E1: every later stage is accepted. E2
  (\op{O5} only): the planted issue survives into the final report. E3: the
  share of the clean run's final-report claims with no content-word partner in
  the condition's report, read against the clean-versus-clean noise floor. This
  endpoint set was chosen before the pass from three candidates defined
  together; the third, an LLM judge, was out of scope.
\item[Overwrite counterfactual.] Production resolves a divergent frequency by
  overwriting it with the recomputed value, at zero calls; it never retries.
  This cell routes the divergence to a retry instead, feeding \arm{A1} the
  supervisor's own correction flags. That measures what the overwrite bought.
  Both readings were fixed before the freeze: if the retry repairs at about
  100\% without collateral damage, the overwrite was a cost decision, not a
  correctness one.
\end{description}

% ---------------------------------------------------------------------
\subsection{Execution}
\label{sec:design-execution}

All trials used \model{} through Google's OpenAI-compatible endpoint, at
temperature~0 except in the warm cell. Fallback to a second model was disabled
throughout. \CSEVEN{The second-model cell used [model, parameter count,
quantisation, digest] served locally by [runtime, version] with a
[32{,}768]-token context, greedy decoding, on [GPU]; its gates and exclusions
are fixed in amendment A10.}

\paragraph{One session per defect}
A temperature of zero does not make a hosted model deterministic across
sessions. Served weights and the serving stack change, and repeated greedy
runs of the same prompt are known to differ~\citep{arimbur2026}. Every
within-defect contrast is therefore matched in time. All arms and repetitions
of one defect run contiguously in a single session. The runner records the
session and refuses to proceed if the iteration order is ever rewritten
arm-outer. For the same reason, no exploratory cell borrows the core's trials
as its comparators. Each re-runs its own comparison condition in-session
(Table~\ref{tab:cells}). Those re-runs double as test--retest measurements of
serving drift, reported in Section~\ref{sec:results-rq2}.

\paragraph{Verification}
Each pass ended with a verification gate. The gate required a complete
checkpoint, one session per defect, and an offline re-scoring of every trial
from its stored output. The re-scoring agreed with the live score in every
trial of every pass. The analysis code was written after the core finished. It
transcribes rules the pre-registration had fixed, was tested on synthetic rows
before it read real data, and writes its own hash into every output. Every
number in Section~\ref{sec:results} is regenerated from the committed trial
records by the committed analysis scripts, with no model calls.

\paragraph{Losses}
The core lost 18 of 1{,}167 trials to provider errors: network drops and one
overnight sleep of the client machine. These cost three defects their
\arm{A1}/\arm{A2} pair. Because all arms of a defect share one session, a loss
cannot depend on the arm. Across all cells the study consumed about 53 million
tokens.

% ---------------------------------------------------------------------
\subsection{Deviations from the pre-registration}
\label{sec:design-deviations}

\begin{table}[tbp]
\centering
\caption{Amendments filed after the freeze. None adds a confirmatory claim or
changes the core. Each was committed before the pass it governs. The last
column records what had already been observed when it was written.}
\label{tab:amendments}
\footnotesize
\begin{tabularx}{\linewidth}{@{}l l Y L{2.9cm}@{}}
\toprule
\# & Date & Change & Results known at filing \\
\midrule
A1 & 2026-09-20 & C1 prompt written; condition factor wired & none \\
A2 & 2026-10-07 & Overwrite counterfactual wired: trigger gate inverted for
  \op{O4}; \arm{A1} fed production's correction flags & core \\
A3 & 2026-10-07 & C1 run as its own pass with in-session C0 comparators & core \\
A4 & 2026-10-07 & Confirmatory analysis code written; one interim budget look
  disclosed & core \\
A5 & 2026-10-07 & Placement pass: all four levels, \arm{P1} re-run in-session
  & core, D19, C1 \\
A6 & 2026-10-07 & \op{O1} severity: level~2 defined (extent), \arm{A1} only &
  core, D19, C1 \\
A7 & 2026-10-07 & Warm-resampling cell, both readings fixed & core, D19, C1,
  placement \\
A8 & 2026-10-08 & Real-defect replay wired; descriptive only & all of the above
  \\
A9 & 2026-10-08 & Downstream propagation brought into scope; endpoint set
  chosen before the pass & all of the above \\
\CSEVEN{A10} & \CSEVEN{date} & \CSEVEN{second-model cell: model identity,
  gates, cells, readings} & \CSEVEN{all Gemini cells} \\
\bottomrule
\end{tabularx}
\end{table}

Table~\ref{tab:amendments} lists every change since the freeze. Three
disclosures go beyond the table.

\emph{An interim look.} On 5~October~2026, after 12 of the 142 core defects,
outcome counts and per-trial token use were inspected to check the budget.
Nothing about the design, the stopping point or the analysis changed, and the
run continued to completion.

\emph{Amendments written with results in hand.} Amendments A5--A9 were filed
after the core's ceiling was known. Their sizes and arms follow a grid fixed in
the study plan before any result existed. The one parameter that grid left
open, \op{O1}'s second severity level, was chosen between two candidates
defined before either ran. The warm cell (A7) is the exception. It was designed
in response to the core's byte-identity pattern, and it is reported as a test
of that explanation, not as an independent finding.

\emph{A re-roled arm.} Before the freeze, \arm{A3} was planned as a near-zero
floor, with two further confirmatory tests built on that assumption: Cochran's
$Q$ across arms and \arm{A2} against \arm{A3}. A pre-freeze measurement showed
that a fresh draw at \arm{P0} regenerates a clean output and was expected to
score about 53\%, against the 5\% the power calculation had assumed. \arm{A3}
was re-roled as a regeneration baseline. Both tests were moved to the secondary
set, and a pre-committed negative result built on \arm{A3} was withdrawn. All
of this happened before the freeze and before any trial, and the
pre-registration records it.
